\section{Practical implications}\label{sec:practical}

\subsection{Sampling frequency for the convexity adjustment}

A variance swap strike converts to a volatility swap strike through a convexity
adjustment that depends on $\operatorname{Var}(\log IV)$. Under lognormal $V$ the
relation is exact and requires no expansion,
\begin{equation}\label{eq:kvol}
K_{\mathrm{vol}} \;=\; \mathbb{E}\!\left[\sqrt{V}\right] \;=\; \sqrt{\mathbb{E}[V]}\,
\exp\!\left(-s^{2}/8\right), \qquad s^{2} = \operatorname{Var}(\log V).
\end{equation}
The second-order expansion in $\operatorname{Var}(V)/\mathbb{E}[V]^{2}$ of
\citet{brockhaus2000}, commonly used in its place, departs from
equation~\eqref{eq:kvol} by ten percent at $\kappa = 0.182$. The measured $\kappa$ % numbers.csv: convex_kappa_boundary
in this sample runs from $1.26$ to $6.84$, and at the upper end of that range the % numbers.csv: convex_kappa_measured_lo, convex_kappa_measured_hi
expansion returns a negative strike, so we use equation~\eqref{eq:kvol}
throughout.

Substituting $\operatorname{Var}(\log RV_M)$ for $\operatorname{Var}(\log IV)$
overstates the adjustment by between $5.9$ and $134.3$ percent at % numbers.csv: convex_overstatement[NQ/GLOBEX/1day], convex_overstatement[ES/RTH/30min]
five-minute-equivalent sampling. Since the overstatement is a function of the
sampling frequency the practitioner chooses, Table~\ref{tab:freq} states, for
each cell, the frequency at which the overstatement falls below five percent.

\begin{table}[htbp]\centering\footnotesize
\caption{Sampling required for the convexity bias to fall below five percent.}
\label{tab:freq}
\begin{tabular}{lcc}
\toprule
Cell & Bias at five-minute sampling & Sub-bar length required \\
\midrule
ES/GLOBEX 1day & $19.3\%$ & $1$ minute \\ % numbers.csv: convex_overstatement[ES/GLOBEX/1day], freq_required_min[ES/GLOBEX/1day]
NQ/GLOBEX 1day & $5.9\%$  & $1$ minute \\ % numbers.csv: convex_overstatement[NQ/GLOBEX/1day], freq_required_min[NQ/GLOBEX/1day]
ES/RTH 1day    & $17.6\%$ & $1$ minute \\ % numbers.csv: convex_overstatement[ES/RTH/1day], freq_required_min[ES/RTH/1day]
NQ/RTH 1day    & $6.9\%$  & $1$ minute \\ % numbers.csv: convex_overstatement[NQ/RTH/1day], freq_required_min[NQ/RTH/1day]
NQ/RTH 1h      & $21.0\%$ & $1$ minute \\ % numbers.csv: convex_overstatement[NQ/RTH/1h], freq_required_min[NQ/RTH/1h]
ES/RTH 1h      & $62.6\%$ & unreachable \\ % numbers.csv: convex_overstatement[ES/RTH/1h], freq_required_min[ES/RTH/1h] = MISSING
ES/RTH 30min   & $134.3\%$ & unreachable \\ % numbers.csv: convex_overstatement[ES/RTH/30min], freq_required_min[ES/RTH/30min] = MISSING
NQ/RTH 30min   & $42.1\%$ & unreachable \\ % numbers.csv: convex_overstatement[NQ/RTH/30min], freq_required_min[NQ/RTH/30min] = MISSING
\bottomrule
\end{tabular}
\end{table}

One-minute sampling suffices at the daily horizon on both instruments and under
both session geometries, but nowhere else. At the thirty-minute horizon on either
instrument, and at the hourly horizon on ES, no frequency available in the grid
brings the bias under five percent. A thirty-minute window does not contain
enough returns for the proxy's own dispersion to fall below the quantity being
measured. A practitioner pricing a volatility swap on a thirty-minute horizon
cannot remedy the bias by sampling faster.

No options data is held, and we report a bias in the adjustment term without any
claim about executable profit.

\subsection{Misclassification at a median threshold}

Where a decision turns on a threshold placed inside the distribution rather than
in its tail, the misclassification rate follows from reliability alone. With
thresholds at the respective medians and lognormal noise, the correlation between
$\log RV_M$ and $\log IV$ is $\sqrt{\lambda}$, and the bivariate normal orthant
probability gives
\begin{equation}\label{eq:mis}
\Pr\left[\text{disagree}\right] \;=\; \frac{\arccos\!\sqrt{\lambda}}{\pi}.
\end{equation}
Equation~\eqref{eq:mis} is exact under those assumptions and is plotted in
Figure~\ref{fig:mis}. A practitioner who has measured $\lambda$ for their own
series can read the implied rate off the figure directly.

The measured cells sit between $7.9$ and $13.6$ percent, while the empirical rate % numbers.csv: mis_analytic[NQ/GLOBEX/insample], mis_analytic[ES/GLOBEX/insample]
measured against a finest-grid noise-robust proxy is lower, falling between $50$
and $77$ percent below the analytic value. The two series are computed from the
same window and their errors are positively correlated, so the empirical figure
should be read as a lower bound. Moving the classification boundary from the
median to the two tercile boundaries raises the empirical rate to between $6.0$ % numbers.csv: mis_tercile_min, mis_tercile_max
and $11.9$ percent. The mechanism there is the count of boundaries, not the
density at them, since two boundaries admit more disagreements than one. For a
unimodal distribution the tercile boundaries sit further from the mode than the
median does.

\subsection{Volatility-target sizing}

We recommend against correcting volatility-target sizing for measured
reliability. On the holdout, replacing the textbook shrinkage weight with the
measured one changes tracking error by at most $1.008$ percent in relative terms,
in the same direction in every cell and at every point of a cost sweep spanning
an eightfold range of transaction cost. Reliability at five-minute-equivalent
sampling runs from $0.40$ to $0.94$ against a textbook $0.67$ to $0.99$, so the % numbers.csv: lambda[ES/RTH/30min/extended], lambda[NQ/GLOBEX/1day/extended], lambda_theory[ES/RTH/30min/extended], lambda_theory[NQ/GLOBEX/1day/extended]
gap being corrected is real and large, but it does not move a position far enough
to matter at a ten percent volatility target with daily rebalancing. The same holds
for inverse-volatility risk parity at monthly rebalancing, where averaging the
estimate over twenty-one sessions divides the bias term by twenty-one.
